\subsection{Acts and agencies}
\label{sec:9.4.2}
Read against the positions of \S\ref{sec:9.3}, the Melt ICE Act is the nearest thing on the record to the removal floor enacted as statute. It closes positions no vendor term could reach, and Table 9.3 shows which.

\runin{Table 9.3. The Melt ICE Act read against the positions where a refusal can be defeated}

\begingroup\small\renewcommand{\arraystretch}{1.12}
\begin{longtable}{@{}>{\raggedright\arraybackslash}p{\dimexpr 0.104\linewidth-2\tabcolsep\relax}>{\raggedright\arraybackslash}p{\dimexpr 0.321\linewidth-2\tabcolsep\relax}>{\raggedright\arraybackslash}p{\dimexpr 0.574\linewidth-2\tabcolsep\relax}@{}}
\hline
\textbf{Position (\S\ref{sec:9.3})} & \textbf{What the Melt ICE Act closes} & \textbf{What it leaves open} \\
\hline
\endhead
8. Handoff & Ends detention and monitoring contracts, then bars all federal funds for detention and monitoring programs, whoever runs them. & The handoff to transport and removal, where it is paid for outside ICE's operations account, and the handoff to a foreign government's custody. \\ \hline
2. Routing & The state and local route: repeals 287(g) and the information-sharing mandates, and defunds DHS partnerships with local police. & The federal route. The bar on civil enforcement applies to one ICE account. Border Patrol, whose agents Amnesty documented in the same operations \autocite{Amnesty2026Whistles}, is not reached, and neither is money appropriated outside that account. Identification by face in the street is not among its provisions. \\ \hline
3. Case description & Bars its own grantees from passing personal information to any federal entity. & Most federal data: the Medicaid and IRS transfers, commercial brokers, the case-management, ImmigrationOS and ELITE systems, and the skip tracers. The bill's “monitoring” means monitoring programs, not analytics. \\ \hline
5–6. Aggregation and review & Ends most detention, and with it the risk tool's recommendation. & Expedited removal, designations under a proclamation, and an adjudicator answerable to the Attorney General. \\ \hline
9. Runtime intervention & — & The halt. Nothing in the bill makes a removal wait for a contest to be decided. \\ \hline
\end{longtable}
\endgroup

The right-hand column has one pattern. Where the bill is written over an act (“no Federal funds may be used with respect to immigration detention and monitoring programs”), it reaches every route, because the act is the same whoever performs it and whichever account pays. Where it is written over an agency or an account (the funds provided to ICE for operations and support), the act can move next door. The two biometric drafts make the same choice in miniature. The first is written over two agencies and leaves every other federal agency free to identify people by their faces. The second is written over the act, by every federal entity, and reaches state and local police through their grants. This is \S\ref{sec:4.1}'s rule, that the floor is stated over acts, applied to a statute, and Congress has learned it at cost before.

The 1984 Boland Amendment barred funds available to the CIA, the Defense Department “or any other agency or entity of the United States involved in intelligence activities” from supporting military operations in Nicaragua \autocite[§8066(a)]{boland1984}. The National Security Council staff ran the program anyway. For sixteen months it used money from third countries, private donors and arms sales to Iran to carry out what the investigating committees called a program Congress believed it had prohibited \autocite{irancontra1987}. Detection came from outside the operation in two ways: a supply plane was shot down over Nicaragua on 5 October 1986, and investigators found a memorandum that had escaped the shredder \autocite{irancontra1987}. The committees' remedy was to write the next rule over the act. They recommended that any department, agency or entity of the government need a presidential finding before a covert action, whatever the source of funds \autocite{irancontra1987}.

The removal pipeline has already shown where its acts move. In the eighteen months to July 2026, ICE held people in three places outside its ordinary contracts \autocite{gao2026detention}:

\begin{itemize}
\item an Army base, under a contract the Army awarded;
\item eight Bureau of Prisons facilities, under an interagency agreement whose rates ICE could not negotiate;
\item a state-run facility for which ICE had no contract at all. DHS reimbursed the state through a FEMA grant program built from funds appropriated for sheltering migrants, and ICE cited a county's 287(g) agreement as its legal basis.
\end{itemize}

Over the same period the number of facilities ICE was authorized to use roughly doubled, from 134 to 272 \autocite{gao2026detention}. A bar written over ICE's detention contracts would have caught none of these three arrangements as they began. The Melt ICE Act's bar on all federal funds for detention would have caught all three. The record shows a fourth place people were held: a foreign government's prison, under the paid arrangement with which this chapter opened \autocite{Renteria2025ElSalvador}. Whether either bar would have caught it is doubtful, since a person held abroad after removal is arguably no longer in immigration detention. A bar written over the act of holding people the United States has removed, wherever and by whomever, would have caught it.
